\phantomsection
\chapter*{Önsöz}
\markboth{Önsöz}{Önsöz}
\addcontentsline{toc}{chapter}{Önsöz}

Bu kitap, lisans düzeyinde uygulamalı istatistik öğrenen okuyucu için
hazırlanmıştır. Temel amacı, araştırma sorusundan savunulabilir bir sonuca
ilerlemeyi öğretmektir: Bu veriyle hangi soru yanıtlanabilir, hangi analiz
neden seçilir ve sonuçtan ne ileri sürülebilir? Formüller ve hesaplamalar,
belirsizliği hesaba katan ve yorumun sınırlarını açıklayan bu düşünme
sürecinin araçlarıdır.
Bu yaklaşım, istatistik eğitiminde araştırma sorusu, veri üretimi, benzetim,
teknoloji ve bağlamsal yorumun birlikte ele alınmasını öneren GAISE
ilkeleriyle uyumludur \citep{gaise2016}.

\ifimoders
Bu ders sürümü, kitabın IMO301 İstatistik ders izlencesine uyarlanmış
14 bölümlük okuma rotasıdır. Haftalık iki saatlik dersin ana hattı;
örnekleme dağılımları, Merkezi Limit Teoremi, nokta ve aralık tahmini,
hipotez testleri, \tstat testleri, kategorik veri analizi, korelasyon ve basit
doğrusal regresyondur. Eğitim verileri üzerindeki uygulamalar hesaplama ile
bağlamsal yorumu birleştirir; son bölüm genel sınava hazırlığa ayrılmıştır.
\else
Kapsamlı sürüm 17 bölümden oluşur. İstatistiksel düşünme ve veri üretiminden
başlayarak örnekleme değişkenliği, tahmin, test ve modellemeye ilerler;
bootstrap ve rastgeleleştirme, çoklu regresyon ve varyans analiziyle kapsamı
genişletir. Son bölüm, gerçek bir eğitim verisi üzerinde yöntem seçimi ve
raporlamayı birleştirir. Bu rota, bütün konuların haftalık iki saatlik tek
bir derste tamamlanması gerektiği anlamına gelmez.
\fi

Kuramsal açıklamalar çözümlü örnekler, sık yapılan hatalar, bölüm sonu
soruları ve araştırma problemleriyle birlikte ele alınmıştır. Yöntemlerin
yalnızca nasıl uygulandığı değil, hangi koşullarda kullanılabileceği ve elde
edilen sonuçların nasıl raporlanması gerektiği de açıklanmıştır. Ortak eğitim
verisi üzerindeki uygulamalar, aynı kayıtlardan hareketle ortalama, oran ve
ilişki hakkında farklı sorular sorulabileceğini gösterir. Bu tekrarlar,
bağımsız araştırmalar değil, sorular ve yorumlar arasındaki farkları
inceleme fırsatıdır.

SPSS, R ve Python uygulamaları, yazılım uzmanlığını başlı başına amaçlamak
yerine analiz tercihlerinin ve çıktıların izlenmesini destekler. Okuyucudan
üç yazılımı birden öğrenmesi beklenmez. Hesaplamaların uyuşması, kullanılan
varsayımların veya bilimsel yorumun doğruluğunun kanıtı değildir; hangi
sonucun nasıl kontrol edildiği uygulamaların açıklamalarında belirtilir.

Metin temel olasılık, cebir ve grafik okuma bilgisi üzerine kuruludur.
Temel istatistiğini yenilemek isteyen başlangıç düzeyi lisansüstü okuyucu
da bu rotadan yararlanabilir. Kitap ileri matematiksel istatistik veya
eksiksiz bir yazılım başvuru kılavuzu olmayı amaçlamaz.
Matematiksel ifadeler kavramsal anlamlarıyla birlikte verilmiş;
hesaplama ile yorum arasında bağ kurulmasına öncelik tanınmıştır. Okuyucular
bölümleri haftalık izlence sırasıyla izleyebileceği gibi formül, dağılım
tablosu ve yazılım eklerini bağımsız bir başvuru kaynağı olarak da
kullanabilir.

Kitaba eşlik eden dijital uygulamalara \url{https://ibrahimguney.github.io/istatistik-lab/} adresinden erişilebilir. Bu kaynağın öğrencilerin, araştırmacıların ve uygulayıcıların istatistiksel düşünme ve modelleme becerilerine katkı sağlaması amaçlanmaktadır.

Bu kitabın hazırlanmasında metinlerin ve LaTeX dosyalarının düzenlenmesi,
kod ve görsellerin geliştirilmesi ile çapraz başvuruların oluşturulması için
Prism'in yapay zekâ destekli araçlarından yararlanılmış olup bilimsel içerik
ve nihai metne ilişkin sorumluluk yazara aittir \citep{openai2026prism}.

\begin{flushright}
İbrahim Güney\\
İstanbul, 2026
\end{flushright}